%!TEX root = TheoreticalPhysicsC.tex


\section{Lecture: Spreading of the Wave Packet}
The term ``wave packet'' is used for spatially confined functions, such as $\Psi\in L^2(\mathbb{R}^3)$, and not for plane waves. From the standpoint of the statistical interpretation there are therefore only such ``packets'', and the discussion of plane waves is at best a slightly abbreviated way of talking about the Fourier transform. In the early days of quantum mechanics, well-localised wave packets were moreover regarded as a kind of substitute for particles. Judged by that
standard, one property of the wave packets proved rather troublesome: they ``flowed apart'', which for particles is really quite unseemly behaviour. The spreading of the wave packet was thereafter seen as a typically quantum mechanical effect. If, however, you read $\Psi$ statistically throughout, there is nothing paradoxical about this process: why, after all, should a probability distribution
not grow broader with time? We shall see here that the spreading of wave packets is precisely a classical effect: asymptotically it is determined exactly by the momentum distribution, just as for a swarm of classical particles.

\begin{greybox}[Why real particles nevertheless do not spread]
	A particle is rarely alone. Matter consists of charged constituents, and the electromagnetic field asks them incessantly ``are you there?'' -- the particle answers by producing a little field, and this ceaseless position measurement localises it again and again. Keyword: \emph{decoherence}.
\end{greybox}

Let us therefore now consider particles in $\mathbb{R}^3$ with Hamiltonian
\begin{equation}
	H = \frac{\widehat{p}_\text{tot.}^2}{2m} = \frac{\widehat{p}_{x}^2}{2m}+\frac{\widehat{p}_{y}^2}{2m}+\frac{\widehat{p}_{z}^2}{2m}.
\end{equation}
In this case $H$ corresponds, as usual in QM, to the energy of a free particle. 

We are interested in the probability distribution of position at various times, especially at large times, that is, in the results of
position measurements when a long time $t$ has elapsed between preparation and measurement. We have two equivalent alternatives for tackling this problem: the Heisenberg picture and the Schrödinger picture. In the Schrödinger picture we would solve the Schrödinger equation 
\begin{equation}
	i\hbar \frac{\partial}{\partial t}\Psi = \widehat{H}\Psi
\end{equation}
and then read off the result as 
\begin{equation}
	p_t(x) \equiv |\Psi(x,t)|^2
\end{equation}
-- a cloud in space, gathered up detector click by detector click. The difference between the pictures is pure bookkeeping: there the state moves, here the measurement, and because $\tr(\rho_t\widehat{F}) = \tr(\rho\,\widehat{F}(t))$, both deliver the same numbers. We shall carry out the calculation in the Heisenberg picture, i.e. we hold $\Psi$ (and with it
the density matrix $\rho$) fixed and compute the operators
\begin{equation}
	\widehat{x}(t) \equiv e^{\frac{it}{\hbar} \widehat{H}}\widehat{x} e^{-\frac{it}{\hbar} \widehat{H}}, \quad \widehat{y}(t) \equiv e^{\frac{it}{\hbar} \widehat{H}}\widehat{y} e^{-\frac{it}{\hbar} \widehat{H}} \quad \text{and} \quad \widehat{y}(t) \equiv e^{\frac{it}{\hbar} \widehat{H}}\widehat{y} e^{-\frac{it}{\hbar} \widehat{H}}.
\end{equation}

Next we look for suitable ``coordinates'' in which we can carry out our calculation. Because the exponential function appears, it is natural to choose a representation of the Hilbert space in which $H$ is a multiplication operator, i.e. the momentum representation:
\begin{equation}
	\Phi(\mathbf{p})\equiv \Phi(p_x,p_y,p_z) \equiv \frac{1}{\sqrt{2\pi\hbar}} \iiint_{-\infty}^\infty \Psi(x)e^{-\frac{i}{\hbar}(p_x x+ p_y y+ p_z z)}\,dxdydz.	
\end{equation} 
Attached to each $\mathbf{p}$ is the plane wave $|\Psi_{\mathbf{p}}\rangle$ from Lecture 8, normalised here so that $\langle\Psi_{\mathbf{p}}|\Psi_{\mathbf{p}'}\rangle = \delta^{(3)}(\mathbf{p}-\mathbf{p}')$:
\begin{equation}
	\langle \mathbf{x}|\Psi_{\mathbf{p}}\rangle = \frac{1}{(2\pi\hbar)^{3/2}}e^{\frac{i}{\hbar}(p_x x + p_y y + p_z z)},
\end{equation}
and these improper kets furnish a resolution of the identity:
\begin{equation}\label{eq:eins-impuls}
	\mathbb{I} = \iiint_{-\infty}^\infty |\Psi_{\mathbf{p}}\rangle\langle\Psi_{\mathbf{p}}|\,dp_xdp_ydp_z.
\end{equation}
The coefficients $\langle\Psi_{\mathbf{p}}|\Psi\rangle$ in this basis are precisely $\Phi(\mathbf{p})$.\footnote{The prefactor above in $\Phi(\mathbf{p})$ is the one-dimensional one; in three dimensions it should be $(2\pi\hbar)^{-3/2}$.} Inserting an identity like this is our most important computational tool; we shall need it again in a moment.

This basis is chosen because the plane waves are eigenvectors of $\widehat{H}$,
\begin{equation}\label{eq:heigen}
	\widehat{H}|\Psi_{\mathbf{p}}\rangle = \frac{|\mathbf{p}|^2}{2m}|\Psi_{\mathbf{p}}\rangle,
\end{equation}
so that in the power series every $\widehat{H}$ turns into a number:
\begin{equation}
	e^{-\frac{it}{\hbar}\widehat{H}}|\Psi_{\mathbf{p}}\rangle = \sum_{n=0}^{\infty}\frac{1}{n!}\left(-\frac{it}{\hbar}\right)^n \widehat{H}^n|\Psi_{\mathbf{p}}\rangle = e^{-\frac{it}{\hbar}\frac{|\mathbf{p}|^2}{2m}}|\Psi_{\mathbf{p}}\rangle.
\end{equation}
In the exponent there was an operator; now there is a number -- hence our obsession with eigenvectors. This is permitted \emph{only} on an eigenvector: replacing a matrix by a number is something you have to earn.

We shall also need in a moment that $\widehat{x}$ acts in the momentum representation as $i\hbar\frac{\partial}{\partial p_x}$. To see this in one dimension you insert the position resolution of the identity,
\begin{equation}
	\langle \Psi_p|\widehat{x}|\Psi\rangle = \int_{-\infty}^\infty \langle \Psi_p|x\rangle\, x\, \Psi(x)\,dx = \frac{1}{\sqrt{2\pi\hbar}}\int_{-\infty}^\infty x\, e^{-\frac{i}{\hbar}px}\,\Psi(x)\,dx,
\end{equation}
produce the factor $x$ as a derivative with respect to $p$,
\begin{equation}
	i\hbar\frac{\partial}{\partial p}e^{-\frac{i}{\hbar}px} = i\hbar\left(-\frac{ix}{\hbar}\right)e^{-\frac{i}{\hbar}px} = x\,e^{-\frac{i}{\hbar}px},
\end{equation}
and pull it out in front of the integral:
\begin{equation}\label{eq:xinp}
	\langle \Psi_p|\widehat{x}|\Psi\rangle = i\hbar\frac{\partial}{\partial p}\Phi(p).
\end{equation}
For $\widehat{y}$ and $\widehat{z}$ it goes exactly the same way -- a nice symmetry with the position representation, where conversely $\widehat{p}$ acts as $-i\hbar\frac{\partial}{\partial x}$: the same structure with the roles interchanged, up to the sign.

Then
\begin{equation}\label{eq:timedepx}
	\begin{split}
		(\widehat{x}(t)\Phi)(\mathbf{p}) &= e^{\frac{it}{\hbar} \frac{|\mathbf{p}|^2}{2m}}\left(\widehat{x}(t)e^{-\frac{it}{\hbar}\widehat{H}}\Phi\right)(\mathbf{p}) \\
		&= e^{\frac{it}{\hbar} \frac{|\mathbf{p}|^2}{2m}}\left(i\hbar \frac{\partial}{\partial p_x}\right)\left(e^{-\frac{it}{\hbar}\frac{|\mathbf{p}|^2}{2m}}\Phi(\mathbf{p})\right) \\
		&= t \frac{p_x}{m}\Phi(\mathbf{p}) + i\hbar \frac{\partial}{\partial p_x} \Phi(\mathbf{p}), \\ 
	\end{split}
\end{equation}
hence
\begin{equation}
	\widehat{x}(t) = t \frac{\widehat{p}_x}{m} + \widehat{x}(0),
\end{equation}
where $\frac{\widehat{p}_x}{m}$ represents the \emph{velocity operator} (in the $x$-direction). Similar expressions hold for $\widehat{y}(t)$ and $\widehat{z}(t)$. The expression $\frac{p_\mu}{m} = \frac{\partial}{\partial p_\mu} \frac{|\mathbf{p}|^2}{2m}$, $\mu = x,y,z$, is also called the \emph{group velocity}, because it describes the velocity of \emph{groups}, that is, of wave packets. The calculation is laborious; what you should remember is its result, not the sequence of steps.

To obtain the time dependence of the expectation value of $\widehat{x}(t)$, all you now have to do is take the trace with $\rho =  |\Psi\rangle\langle\Psi|$ on both sides of (\ref{eq:timedepx}). We thereby obtain
\begin{equation}
	\begin{split}
		\langle\widehat{x}\rangle(t) &= \tr(\rho \widehat{x}(t)) \\
		&= \tr\left(\rho \left(t \frac{\widehat{p}_x}{m} + \widehat{x}(0)\right)\right)\\
		&= t \frac{\langle\widehat{p}_x \rangle(0)}{m} + \langle\widehat{x}\rangle(0).
	\end{split}
\end{equation}
That is exactly the newtonian rule ``position equals initial position plus velocity times time''.

\begin{intuition}
	The same formula, different physics: classically a particle \emph{has} a position, and expectation values are not needed at all. Here only the mean runs classically, while the width grows -- a classical cloud with sharp position and momentum does not spread, a quantum mechanical one must. Classical mechanics is the limiting case $\hbar\to 0$.
\end{intuition}

The operator equation contains far more information than just that, however: the diagonalisation of the operators $\widehat{x}(t)$ does after all contain all the information about the probability distributions of position at time $t$ in arbitrary states. The appearance of the position probability distribution at large times is therefore completely determined by the joint distribution of the operators $\frac{\widehat{p}_\mu}{m}$. The behaviour of $\widehat{x}$, $\widehat{y}$ and $\widehat{z}$ therefore shows no difference whatever between classical and quantum mechanical particles.

% fig_vl10_zerfliessen.tex — Zerfliessen des Wellenpakets: the packet
% drifts uniformly while it spreads, and its width settles onto exactly
% the straight line a swarm of classical particles with the same
% momentum spread would follow.
%
% GEOMETRY FROM THE FORMULA.  Everything here follows from the operator
% identity this lecture derives,
%
%       x(t) = x(0) + (t/m) p ,
%
% by taking mean and variance in a state with no initial x-p
% correlation:
%
%       <x>(t)  = <x>(0) + <p> t/m ,
%       Dx(t)^2 = Dx(0)^2 + (Dp t/m)^2 .
%
% With tau = m Dx(0)/Dp — and for a minimum-uncertainty packet,
% Dx(0) Dp = hbar/2, hence tau = 2 m Dx(0)^2/hbar — the second line is
%
%       Dx(t) = Dx(0) sqrt(1 + (t/tau)^2),
%
% the hyperbola of panel (b).  Its asymptote Dp t/m is precisely how a
% classical swarm with momentum spread Dp comes apart, so the claim of
% the lecture — that the spreading is asymptotically fixed by the
% momentum distribution alone — is the statement that the two curves in
% (b) merge.  Panel (a) plots the Gaussian carrying those two moments,
%
%       |Psi(x,t)|^2 = exp(-(x-<x>(t))^2/(2 Dx(t)^2)) / (sqrt(2 pi) Dx(t)),
%
% so every centre, width and height is forced by the formula rather than
% drawn by eye.  Lengths in units of Dx(0), times in units of tau, drift
% fixed by <p> = 3 Dp: the centres then sit at 0, 3, 6, 9 and the widths
% at 1, sqrt2, sqrt5, sqrt10.
%
% DESIGN DECISIONS
%   * colour encodes TIME and nothing else — entries 1, 2, 4, 5 of the
%     ordered `tpclevels` ramp (navy, blue, orange, red), the same four
%     in both panels, so a curve in (a) can be found again as a dot in
%     (b).  Entry 3 (munsell) is skipped on purpose: next to cgblue it
%     is the same teal to the eye, and here the two would be neighbours
%     in time rather than well separated as in fig_vl11;
%   * the width curve in (b) is green because it is a different
%     QUANTITY, not a fifth time; its formula is set in the same green,
%     which spares the panel a legend;
%   * grey is reference furniture only, and every grey annotation is
%     set in grey too: the verticals at <x>(t), whose EQUAL SPACING is
%     uniform motion; the dashed envelope through the peaks, which
%     falls as 1/Dx(t) exactly as fast as the widths grow, so the
%     conserved area is shown rather than asserted, and whose label is
%     set in the empty strip immediately above it rather than in the
%     right margin, where it crowded the t=3tau label; the initial width
%     and the classical asymptote in (b);
%   * time labels sit to the RIGHT of each peak rather than above it,
%     which is the only side on which the descending envelope leaves
%     room;
%   * panel (b) uses `scale only axis` with a 4.2 x 5.4 cm box on ranges
%     4.2 x 5.4, i.e. 1 cm per unit on both axes.  The asymptote is then
%     a true 45-degree line and its label can be set along it.  The y
%     range runs to 5.4 although the curve reaches only sqrt(1+4.2^2) =
%     4.32 inside the frame: that headroom is where the three annotation
%     lines sit, so none of them is written over the curve.  The widest
%     of them, the green formula, is 3.8 cm in the fonts of the notes —
%     wider than in the cmfonts preview, which is why the panel is sized
%     against the former;
%   * both panels are 5.4 cm tall so that titles and x axes line up;
%   * the orders of magnitude for tau — 10^-16 s for an electron
%     localised to an angstrom, 10^13 s for a microgram grain — are
%     NOT in the panel.  They were, at its foot, and in the fonts of the
%     notes (wider than the cmfonts preview) the grey asymptote ran
%     straight through them; a 4.2 cm panel crossed by two reference
%     lines has no free region 3.4 cm wide, so the numbers now stand in
%     the sentence below the figure, where the definition of tau in (b)
%     hands them straight over.

\begin{center}
\begin{tikzpicture}

% ─────────────────────────────────────────────────────────────
%  (a)  |Psi(x,t)|^2 at four times
% ─────────────────────────────────────────────────────────────
\begin{axis}[
    name=packets,
    tpcplot,
    scale only axis, width=6.4cm, height=5.4cm,
    grid=none,
    xmin=-3.2, xmax=16.8,
    ymin=0,    ymax=0.47,
    xtick={0,3,6,9,12,15},
    ytick={0,0.1,0.2,0.3,0.4},
    xlabel={$x/\Delta x(0)$},
    ylabel={$\Delta x(0)\,|\Psi(x,t)|^2$},
    every axis x label/.style={at={(axis description cs:0.5,-0.15)},
                               anchor=north, font=\small},
    every axis y label/.style={at={(axis description cs:-0.16,0.5)},
                               rotate=90, anchor=south, font=\small},
    title={(a) \tpcifenglish{The packet spreads}{Das Paket zerfließt}},
    samples=401,
]

  % Centres <x>(t) = 3 t/tau, each rule drawn only up to the peak it
  % belongs to.  Their equal spacing is the newtonian rule.
  \draw[levellinegrey] (axis cs:0,0) -- (axis cs:0,0.3989);
  \draw[levellinegrey] (axis cs:3,0) -- (axis cs:3,0.2821);
  \draw[levellinegrey] (axis cs:6,0) -- (axis cs:6,0.1784);
  \draw[levellinegrey] (axis cs:9,0) -- (axis cs:9,0.1262);

  % Envelope of the peaks: 1/(sqrt(2 pi) Dx(t)) with t/tau = s/3.  It
  % passes exactly through the four maxima.
  \addplot[envelope, domain=0:12.6, samples=140]
      {1/(sqrt(2*pi)*sqrt(1+x^2/9))};

  % The packets.  Widths sqrt(1+(t/tau)^2) = 1, sqrt2, sqrt5, sqrt10.
  % Each runs across the whole panel: a curve stopped short leaves a
  % blunt end lying on the axis, which reads as a cut-off rather than
  % as a tail that has died away.
  \addplot[wavefn, spacecadet,   domain=-3.2:16.8]
      {exp(-(x)^2/2)/sqrt(2*pi)};
  \addplot[wavefn, cgblue,       domain=-3.2:16.8]
      {exp(-(x-3)^2/4)/sqrt(2*pi*2)};
  \addplot[wavefn, birrenorange, domain=-3.2:16.8]
      {exp(-(x-6)^2/10)/sqrt(2*pi*5)};
  \addplot[wavefn, birrenred,    domain=-3.2:16.8]
      {exp(-(x-9)^2/20)/sqrt(2*pi*10)};

  % Direct labels beside the peaks, each in the colour of its curve and
  % set just above the peak height, where the envelope has already
  % dropped away.
  \node[lbl, text=spacecadet,   anchor=south west]
    at (axis cs:0.55,0.410) {$t=0$};
  \node[lbl, text=cgblue,       anchor=south west]
    at (axis cs:3.45,0.293) {$t=\tau$};
  \node[lbl, text=birrenorange, anchor=south west]
    at (axis cs:6.45,0.189) {$t=2\tau$};
  \node[lbl, text=birrenred,    anchor=south west]
    at (axis cs:9.45,0.136) {$t=3\tau$};

  % What the grey verticals say, and what the grey envelope says.
  \node[lbl, text=birrengray, font=\scriptsize, anchor=north east,
        align=right] at (axis cs:16.6,0.468)
    {\tpcifenglish{centres}{Zentren}:
     $\langle\widehat{x}\rangle(t)=\langle\widehat{p}\rangle\,t/m$\\[1pt]
     \tpcifenglish{area}{Fl\"ache}: $\int|\Psi|^2dx=1$};
  \node[lbl, text=birrengray, font=\scriptsize, anchor=south west]
    at (axis cs:4.7,0.232)
    {\tpcifenglish{peak height}{H\"ohe} $\propto 1/\Delta x(t)$};
\end{axis}

% ─────────────────────────────────────────────────────────────
%  (b)  The width: hyperbola against its classical asymptote
% ─────────────────────────────────────────────────────────────
\begin{axis}[
    name=widths,
    tpcplot,
    at={(packets.south east)}, anchor=south west, xshift=17mm,
    scale only axis, width=4.2cm, height=5.4cm,
    xmin=0, xmax=4.2, ymin=0, ymax=5.4,
    xtick={0,1,2,3,4}, ytick={0,1,2,3,4},
    grid=major,
    xlabel={$t/\tau$},
    ylabel={$\Delta x(t)/\Delta x(0)$},
    every axis x label/.style={at={(axis description cs:0.5,-0.15)},
                               anchor=north, font=\small},
    every axis y label/.style={at={(axis description cs:-0.21,0.5)},
                               rotate=90, anchor=south, font=\small},
    title={(b) \tpcifenglish{Width}{Breite} $\Delta x(t)$},
]

  % Reference furniture: the classical asymptote Dp t/m and the width
  % the packet started with.  They cross at t = tau.
  \addplot[levellinegrey, domain=0:4.2, samples=2] {x};
  \draw[levellinegrey] (axis cs:0,1) -- (axis cs:1,1);

  % Dx(t)/Dx(0) = sqrt(1+(t/tau)^2)
  \addplot[birrengreen, line width=1.1pt, domain=0:4.2, samples=140]
      {sqrt(1+x^2)};

  % The four times of panel (a), in their colours.
  \addplot[only marks, mark=*, mark size=2.0pt, spacecadet]
      coordinates {(0,1)};
  \addplot[only marks, mark=*, mark size=2.0pt, cgblue]
      coordinates {(1,1.41421)};
  \addplot[only marks, mark=*, mark size=2.0pt, birrenorange]
      coordinates {(2,2.23607)};
  \addplot[only marks, mark=*, mark size=2.0pt, birrenred]
      coordinates {(3,3.16228)};

  % The curve names itself, in its own colour; the grey lines below it
  % state the asymptotics and fix the time scale.
  \node[text=birrengreen, font=\scriptsize, anchor=north west]
    at (axis cs:0.08,5.32)
    {$\Delta x(t)=\Delta x(0)\sqrt{1+(t/\tau)^{2}}$};
  \node[lbl, text=birrengray, font=\scriptsize, anchor=north west]
    at (axis cs:0.08,4.68)
    {$\longrightarrow\ \Delta p\,t/m$ \tpcifenglish{for}{f\"ur} $t\gg\tau$};
  \node[lbl, text=birrengray, font=\scriptsize, anchor=north west]
    at (axis cs:0.08,4.14) {$\tau=2m\,\Delta x(0)^2/\hbar$};

  % The asymptote, labelled along itself — the panel is drawn at
  % 1 cm per unit on both axes, so the line really does run at 45 deg.
  \node[lbl, text=birrengray, font=\scriptsize, rotate=45,
        anchor=north west, inner sep=1.5pt] at (axis cs:1.65,1.50)
    {\tpcifenglish{classical swarm}{klassischer Schwarm}};
\end{axis}

\end{tikzpicture}
\end{center}


In~(a) the grey marks of the centres advance in equal steps, while
the height falls exactly as far as the width grows. In~(b) notice
how early the curve swings onto the classical straight line: after only a few
$\tau$ the difference can no longer be seen. Whether you ever set eyes
on the spreading at all is decided by the time scale $\tau$: for an electron with
$\Delta x(0) = 10^{-10}$\,m it is $\tau \sim 10^{-16}$\,s, whereas for a grain of dust of
$1\,\mu$g with $\Delta x(0) = 10^{-6}$\,m it is $\tau \sim 10^{13}$\,s, i.e.
roughly $300\,000$ years.

The more interesting differences between classical mechanics and quantum mechanics (and the more important applications of quantum mechanics) are to be found, however, with particles in external potentials, or in interaction with one another.

The next step in that direction is a modest one: a particle in a box of \emph{finite} size -- our first example of what ``quantum'' in quantum mechanics means.
